% SOURCE_AUTHORITY: sga5_fr_workpass.tex, lines 13521--13547 (LF-normalized unit).
% SOURCE_SHA256: 791F4EFFC5E02832D5D77ED03518C8156D6F07E4C8238B03545DB93D883FBB28
É claro que a condição a) implica $f'_R(x) - x \in \fm^r$ qualquer que seja $x \in R$, o que dá sentido à condição 2).

Esta observação permite também definir um homomorfismo de grupos abelianos:
\[
\fm \longrightarrow \fm^r/\fm^{r+1}
\]
\[
x \longmapsto (f'_R(x) - x) \bmod \fm^{r+1}
\]

Este homomorfismo se anula sobre $\fm^2$. De fato, se $x_1, x_2 \in \fm$, então $f'_R(x_1x_2) - x_1x_2 \in \fm^{r+1}$, pois, graças à condição a), tem-se:
\[
f'_R(x_1) = x_1 + \alpha_1 \; , \; \alpha_1 \in \fm^r \quad ,
\]
\[
f'_R(x_2) = x_2 + \alpha_2 \; , \; \alpha_2 \in \fm^r \quad ,
\]
e, consequentemente:
\[
f'_R(x_1x_2) = f'_R(x_1)f'_R(x_2) = (x_1 + \alpha_1)(x_2 + \alpha_2) = x_1x_2 + x_1\alpha_2 + x_2\alpha_1 + \alpha_1\alpha_2 \quad ,
\]
e evidentemente $x_1\alpha_2 + x_2\alpha_1 + \alpha_1\alpha_2 \in \fm^{r+1}$.

Resulta, pois, por passagem ao quociente, um homomorfismo de grupos abelianos:
\[
u(f'_R) : \fm/\fm^2 \longrightarrow \fm^r/\fm^{r+1} \quad .
\]
